\documentclass[docmute]{article}
\input{preamble.tex}

\input{watermark/watermark.tex}


\title{Counting (Combinatorics)}
\author{Tomasz Brengos \\  
Committers : Mykhailo Moroz}
\begin{document}
\maketitle

\section{Rule of Sum (Addition Principle)}
If a set $S$ is partitioned into disjoint subsets,
\[
S = S_1 \cup S_2 \cup \cdots \cup S_k,
\]
then the total number of elements in $S$ is the sum of the number of elements in each subset:
\[
|S| = |S_1| + |S_2| + \cdots + |S_k|.
\]

\paragraph{Example:}  
Suppose we wish to count the number of ways to choose a subset of a set $X$ of size $n$, but we only consider subsets of a fixed size $k$. If we let $S$ be the family of all such subsets, then using the rule of sum by dividing the choices according to a distinguished element (say, whether a chosen element is included or not) we can count the subsets by summing over the possibilities. (This idea is used later in proofs for binomial coefficients and the power set.)

\paragraph{Theorem:}
For any set \( X \) with \( |X| = n \) and integers \( 1 \le k \le n \),
\[
\binom{n}{k} = \binom{n-1}{k-1} + \binom{n-1}{k},
\]
with the convention that \( \binom{m}{j} = 0 \) whenever \( j > m \).

\paragraph{Proof:}
Consider \( S = \binom{X}{k} \), the set of all subsets of \( X \) of size \( k \).  
Take any element \( a \in X \). Define:  
 \( S_1 \) as the subsets in \( S \) that contain \( a \) and  
 \( S_2 \) as the subsets in \( S \) that do not contain \( a \).  

Since every element of \( S \) (i.e.\ every $k$-subset of $X$) either contains \( a \) or does not, we see that \( S_1 \) and \( S_2 \) are disjoint and their union forms \( S \), i.e.,  
\[
S_1 \cup S_2 = S.
\]

By the rule of sum, we get:  
\[
|S| = |S_1| + |S_2|.
\]

Now,  
each subset in \( S_1 \) must contain \( a \), so we choose the remaining \( k-1 \) elements from \( X \setminus \{a\} \), which has \( n-1 \) elements. Thus, \( |S_1| = \binom{n-1}{k-1} \).  
Each subset in \( S_2 \) does not contain \( a \), so we choose all \( k \) elements from \( X \setminus \{a\} \). Thus, \( |S_2| = \binom{n-1}{k} \).  

Therefore,  
\[
\binom{n}{k} = |S| = |S_1| + |S_2| = \binom{n-1}{k-1} + \binom{n-1}{k}.
\]

\paragraph{Example:}
Let \( X = \{ \triangle, \square, \circ \} \) (so \( n = 3 \)) and \( k = 2 \), and let \( S = \binom{X}{2} \) be the family of all 2-element subsets of \( X \). Choose \( a = \circ \).
Fixing \( \circ \) as one of the elements in the subset of size \( k \), we get:
  \[
  S_1 = \{\{\circ, \triangle\}, \{\circ, \square\}\}.
  \]
Taking all subsets of size \( k \) without \( \circ \):
  \[
  S_2 = \{\{\triangle, \square\}\}.
  \]
We have \( |S_1| = 2 \) and \( |S_2| = 1 \), so \( |S| = |S_1| + |S_2| = 3 \).

\noindent On the other hand, the identity predicts
\begin{align*}
\binom{3}{2} = \binom{2}{1} + \binom{2}{2} = 2 + 1 = 3,
\end{align*}
which matches: \( |S_1| = \binom{2}{1} \) and \( |S_2| = \binom{2}{2} \), verifying the identity.

\section{Rule of Product (Multiplication Principle)}
When an object is constructed by a sequence of choices, where:
\begin{itemize}[nosep]
    \item The first choice can be made in $a$ ways,
    \item The second in $b$ ways,
    \item $\ldots$
\end{itemize}
the total number of objects is the product:
\[
a \times b \times \cdots.
\]

\paragraph{Example:}
A \emph{word} of length $n$ (or \emph{$n$-word}) over an alphabet $X$ is a sequence $(x_1, x_2, \dots, x_n)$ of elements of $X$ (repetitions allowed). For instance, a word of length $n$ over the binary alphabet $\{0,1\}$ is formed by choosing one of $2$ possibilities for each position. Hence, there are
\[
2^n
\]
possible binary words.

\section{Rule of Bijection}
If there exists a bijection (a one-to-one and onto mapping) between two sets $S$ and $T$, then they have the same number of elements:
\[
|S| = |T|.
\]

\paragraph{Example:}  
Consider the power set of a set $X$, denoted by $\mathcal{P}(X)$. There is a natural bijection between $\mathcal{P}(X)$ and the set of binary sequences of length $|X|$: for each subset $A \subseteq X$, assign the sequence $(a_1,a_2,\dots,a_n)$ where 
\[
a_i = \begin{cases} 
1, & \text{if } x_i \in A, \\
0, & \text{if } x_i \notin A.
\end{cases}
\]
This shows that 
\[
|\mathcal{P}(X)| = 2^{|X|}.
\]

\section{Counting in Two Ways}  

\paragraph{Rule of Counting in Two Ways}  
When two formulae enumerate the same quantity, they must be equal.  

\paragraph{Example:}  
\[
\sum_{i=1}^{n} i = \frac{n(n+1)}{2}
\]

\paragraph{Proof:}  
Consider a lattice grid of size \( (n+1) \times (n+1) \), defined as:  
\[
X = \{(i, j) \mid i, j \in \{1, 2, \dots, n+1\} \}.
\]  
Clearly, \( |X| = (n+1)^2 \).  

Now, partition \( X \) into three subsets:
\begin{itemize}[nosep]
  \item \( X_1 \), the points strictly below the secondary diagonal.
  \item \( X_2 \), the points strictly above the secondary diagonal.
  \item \( X_3 \), the points on the secondary diagonal itself.
\end{itemize}

Since these three sets form a partition, we have:  
\[
|X| = |X_1| + |X_2| + |X_3|.
\]  
Observing their sizes:  
\[
|X_1| = |X_2| = 1 + 2 + \dots + n, \quad |X_3| = n+1.
\]  
Thus,  
\[
(n+1)^2 = 2(1 + 2 + \dots + n) + (n+1).
\]  
Rearranging, we get:  
\[
1 + 2 + \dots + n = \frac{(n+1)^2 - (n+1)}{2} = \frac{n(n+1)}{2}.
\]

Hence, we have proven the formula:  
\[
\sum_{i=1}^{n} i = \frac{n(n+1)}{2}.
\]


\section{Binomial Coefficients and Permutations}
Let $X$ be a set with $|X| = n$.

\paragraph{Subsets:}  
The number of ways to choose a $k$-subset of $X$ is given by the binomial coefficient
\[
\binom{n}{k}.
\]

\paragraph{Permutations:}  
A \( k \)-permutation of a set \( X \) of size \( n \) is a \( k \)-word over the alphabet \( X \) whose entries are distinct.  

\paragraph{Theorem:}  
There are exactly  
\[
n (n-1) (n-2) \dots (n-k+1)
\]  
\( k \)-permutations of an \( n \)-set.  

\paragraph{Question:}  
How are \( k \)-permutations of an \( n \)-set related to \( k \)-subsets of an \( n \)-set?  

\paragraph{Answer:}  
The difference between a \( k \)-permutation and a \( k \)-subset is that a permutation is ordered, while a subset is not.  
To express a \( k \)-permutation in terms of a \( k \)-subset, we need to account for all possible arrangements of the elements, which is \( k! \).
Thus,
\[
\text{number of \( k \)-permutations} = \binom{n}{k} \cdot k!
\]
Expressing \( \binom{n}{k} \) as
\[
\binom{n}{k} = \frac{n!}{k!(n-k)!}
\]
we obtain:
\[
\text{number of \( k \)-permutations} = \frac{n!}{(n-k)!}.
\]


\paragraph{Proof by Counting in Two Ways:}  
Count the number of $k$-permutations of an $n$-set in two ways:
\begin{enumerate}[label=(\arabic*)]
    \item Directly, by applying the rule of product:
    \[
    n \times (n-1) \times \cdots \times (n-k+1) = \frac{n!}{(n-k)!}.
    \]
    \item First choose a $k$-subset (in $\binom{n}{k}$ ways) and then arrange it (in $k!$ ways), giving
    \[
    \binom{n}{k} \cdot k!.
    \]
\end{enumerate}
Equate these two counts to obtain the relation.

\section{Binomial Theorem}
For any real numbers $x,y$ and nonnegative integer $n$, the binomial theorem states:
\[
(x+y)^n = \sum_{k=0}^{n} \binom{n}{k} x^k y^{n-k}.
\]
\paragraph{Explanation:}  
This theorem is a direct consequence of counting the number of ways to choose $k$ copies of $x$ (and the remaining $n-k$ copies of $y$) when expanding the product.

\section{Multisets}  

\paragraph{Definition:}  
A multiset of a set \( X \) of size \( n \) is a function  
\[
m: X \to \mathbb{N}
\]
that assigns a non-negative integer to each element of \( X \), representing its multiplicity in the multiset.  

\paragraph{Example:}
Let \( X = \{a, b, c\} \), and consider the multiset \( \{a, a, b\} \). Then, the function \( m \) is given by:
\[
m(a) = 2, \quad m(b) = 1, \quad m(c) = 0.
\]

The \emph{size} of a multiset \( m \) is the total number of elements counted with multiplicity, i.e.\ \( \sum_{x \in X} m(x) \). A \emph{\( k \)-multiset} is a multiset of size \( k \); in the example above, \( \{a,a,b\} \) is a \( 3 \)-multiset.

\paragraph{Question:}
What is the number of \( k \)-multisets of a set of size \( n \)?

\paragraph{Theorem:}  
The number of all \( k \)-multisets of an \( n \)-set is  
\[
\binom{n+k-1}{k}.
\]  

\paragraph{Proof:}  
Let \( X \) be the set of all \( k \)-multisets of an \( n \)-set.  
Let \( Y \) be the set of all distributions of \( k \) identical objects into \( n \) buckets.  

\paragraph{Claim 1:}
There is a bijection from \( X \) to \( Y \).
Indeed, write the underlying \( n \)-set as \( \{x_1, \dots, x_n\} \) and think of bucket \( i \) as belonging to element \( x_i \). A \( k \)-multiset \( m \) corresponds to the distribution that puts exactly \( m(x_i) \) objects into bucket \( i \); since \( \sum_i m(x_i) = k \), all \( k \) objects are used. Conversely, every distribution of \( k \) identical objects determines the multiplicities of a unique \( k \)-multiset. Thus, by the rule of bijection, we have
\[
|X| = |Y|.
\]

\paragraph{Claim 2:}
Let \( Z \) be the set of all binary sequences of length \( n+k-1 \) with exactly \( n-1 \) ones (or equivalently, \( k \) zeros).
There is a bijection from \( Y \) to \( Z \) (the ``stars and bars'' encoding).
Given a distribution placing \( k_i \) objects into bucket \( i \) (so \( k_1 + \cdots + k_n = k \)), encode it as the sequence
\[
\underbrace{0 \cdots 0}_{k_1} \, 1 \, \underbrace{0 \cdots 0}_{k_2} \, 1 \, \cdots \, 1 \, \underbrace{0 \cdots 0}_{k_n},
\]
i.e.\ the zeros represent the objects and the \( n-1 \) ones are separators between consecutive buckets. The resulting sequence has length \( k + (n-1) \) with exactly \( n-1 \) ones. Conversely, any such binary sequence is split by its ones into \( n \) (possibly empty) blocks of zeros, recovering a unique distribution. Hence,
\[
|Y| = |Z| \Rightarrow |X| = |Z|.
\]
Since the number of such binary sequences is given by  
\[
\binom{n+k-1}{k},
\]  
we conclude that  
\[
|X| = \binom{n+k-1}{k}.
\]


\section{Lattice Paths}
Consider an $m \times n$ grid with lattice points at the intersections.  
\paragraph{Problem:}  
How many paths are there from $(0,0)$ to $(m,n)$ if one may only move right or up?
\paragraph{Solution:}  
Every path consists of exactly $m$ right moves and $n$ up moves. Thus, a path can be represented as a sequence of $m+n$ moves, where we choose $n$ positions (out of $m+n$) for the up moves. Hence, the number of paths is:
\[
\binom{m+n}{n}.
\]

\paragraph{Bijection Explanation:}  
There is a bijection between the set of such lattice paths and the set of binary sequences of length $m+n$ with exactly $n$ ones (representing the up moves).
\end{document}